\RequirePackage{plautopatch}
\documentclass[14pt,aspectratio=169,xcolor=dvipsnames,table,onlytextwidth,dvipdfmx]{beamer}
\usepackage[utf8]{inputenc}
\usepackage{bxdpx-beamer} % dvipdfmxなので必要

\usetheme{Boadilla}
%Beamerフォント設定
\usefonttheme{professionalfonts}
\usepackage[T1]{fontenc}
\usepackage[deluxe,uplatex]{otf} % 日本語多ウェイト化
\usepackage{mlmodern}  % 太いComputer Modern
% MLmodernのバグを修正: cf. https://tex.stackexchange.com/questions/646333/size-of-integral-symbol-in-section-header-with-mlmodern
\DeclareFontFamily{OMX}{mlmex}{}
\DeclareFontShape{OMX}{mlmex}{m}{n}{%
   <->mlmex10%
   }{}
\usepackage{newtxtext} % 数式以外をTXフォントで上書き
% \usepackage[noalphabet,unicode]{pxchfon}
% \setminchofont{KozMinPro-Regular.otf}% 游明朝Regular
% \setboldminchofont{KozMinPro-Bold.otf}% 游明朝Demibold
% \setgothicfont[0]{BIZ-UDGothicR.ttc}% BIZ UDゴシックR
% \setboldgothicfont[0]{BIZ-UDGothicB.ttc}% BIZ UDゴシックB
\renewcommand{\familydefault}{\sfdefault}  % 英文をサンセリフ体に
\renewcommand{\kanjifamilydefault}{\gtdefault}  % 日本語をゴシック体に
\usefonttheme{structurebold} % タイトル部を太字
\setbeamerfont{alerted text}{series=\bfseries} % Alertを太字
\setbeamerfont{section in toc}{series=\mdseries} % 目次は太字にしない
\setbeamerfont{frametitle}{size=\Large} % フレームタイトル文字サイズ
\setbeamerfont{title}{size=\LARGE} % タイトル文字サイズ
\setbeamerfont{date}{size=\small}  % 日付文字サイズ
\usepackage{bxcoloremoji}

% Babel (日本語の場合のみ・英語の場合は不要)
\uselanguage{japanese}
\languagepath{japanese}
\deftranslation[to=japanese]{Theorem}{定理}
\deftranslation[to=japanese]{Lemma}{補題}
\deftranslation[to=japanese]{Example}{例}
\deftranslation[to=japanese]{Examples}{例}
\deftranslation[to=japanese]{Definition}{定義}
\deftranslation[to=japanese]{Definitions}{定義}
\deftranslation[to=japanese]{Problem}{問題}
\deftranslation[to=japanese]{Solution}{解}
\deftranslation[to=japanese]{Fact}{事実}
\deftranslation[to=japanese]{Proof}{証明}
\deftranslation[to=japanese]{Corollary}{系}
\setbeamertemplate{theorems}[normal font] %日本語の場合，定理内を斜体にしない
% overlay-aware proof environment
\renewenvironment<>{proof}{%
\par
\begin{actionenv}#1%
\structure{(証明)}}
{\qed\end{actionenv}}

%Beamer色設定
\definecolor{UniBlue}{RGB}{0,150,200}
\definecolor{UniGreen}{RGB}{3,175,122}
\definecolor{UniPink}{RGB}{255,128,255}
\definecolor{AlertOrange}{RGB}{255,76,0}
\definecolor{AlmostBlack}{RGB}{38,38,38}
\setbeamercolor{normal text}{fg=AlmostBlack}  % 本文カラー
\setbeamercolor{structure}{fg=UniBlue} % 見出しカラー
\setbeamercolor{block title}{fg=UniBlue!50!black} % ブロック部分タイトルカラー
\setbeamercolor{alerted text}{fg=AlertOrange} % \alert 文字カラー
\setbeamercolor{bibliography entry author}{fg=AlmostBlack} % Bib 文字カラー
\setbeamercolor{bibliography entry note}{fg=AlmostBlack} % Bib 文字カラー
\mode<beamer>{
    \definecolor{BackGroundGray}{RGB}{254,254,254}
    \setbeamercolor{background canvas}{bg=BackGroundGray} % スライドモードのみ背景をわずかにグレーにする
}

%フラットデザイン化
\setbeamertemplate{blocks}[rounded] % Blockの影を消す
\useinnertheme{circles} % 箇条書きをシンプルに
\setbeamertemplate{navigation symbols}{} % ナビゲーションシンボルを消す
\setbeamertemplate{footline}[frame number] % フッターはスライド番号のみ

%タイトルページ
\setbeamertemplate{title page}{%
    \vspace{2.5em}
    {\usebeamerfont{title} \usebeamercolor[fg]{title} \inserttitle \par}
    {\usebeamerfont{subtitle}\usebeamercolor[fg]{subtitle}\insertsubtitle \par}

    \vspace{1.5em}
    \begin{columns}[]
        \begin{column}{.35\linewidth}
        \usebeamerfont{titlegraphic}\inserttitlegraphic\par
        \end{column}
        \begin{column}{.6\linewidth}\setlength\topsep{0pt}
            \begin{flushright}
        \usebeamerfont{author}\insertauthor\par
        \usebeamerfont{institute}\insertinstitute \par
        \vspace{3em}
        \usebeamerfont{date}\insertdate\par
            \end{flushright}
        \end{column}
    \end{columns}
}

%biblatex
% \usepackage[style=ext-authoryear,citestyle=authoryear,natbib=true,giveninits=true,maxbibnames=99,maxcitenames=10,url=false,isbn=false,doi=false,articlein=false]{biblatex}
% \DeclareNameAlias{author}{first-last}
% \addbibresource{shrunk.bib}
% \newcommand{\mkbibbracketscol}[1]{\textcolor{UniBlue}{\scriptsize \mkbibbrackets{#1}}}
% \DeclareCiteCommand{\cite}[\mkbibbracketscol]
%   {\usebibmacro{prenote}}
%   {\usebibmacro{citeindex}%
%    \usebibmacro{cite}}
%   {\multicitedelim}
%   {\usebibmacro{postnote}}
% % \renewcommand{\cite}[1]{\textcolor{UniBlue}{\scriptsize\citep{#1}}}
% \renewcommand*{\finalnamedelim}{\addcomma\addspace} % remove "and" before the last author
% \DeclareDelimFormat{nameyeardelim}{\addspace}
% \setbeamertemplate{bibliography item}{}

% Algorithm系
\usepackage{algorithm}
\usepackage[noend]{algorithmic}
\algsetup{linenosize=\color{fg!50}\footnotesize}
\renewcommand\algorithmicdo{:}
\renewcommand\algorithmicthen{:}
\renewcommand\algorithmicrequire{\textbf{Input:}}
\renewcommand\algorithmicensure{\textbf{Output:}}
\renewcommand\algorithmiccomment[1]{\hfill\textcolor{fg!50}{\footnotesize\textit{// #1}}}

% TikZ
\usepackage{tikz}
\usetikzlibrary{positioning,shapes,arrows,quotes,shapes.callouts,calc,graphs,graphs.standard,fit,backgrounds,perspective,matrix}
\tikzset{point/.style={circle, fill=AlertOrange, inner sep=1pt, text width=1pt}}
\tikzset{bpoint/.style={circle, fill=black, inner sep=1pt, text width=1pt}}
\tikzset{alt/.code args={<#1>#2#3}{\alt<#1>{\pgfkeysalso{#2}}{\pgfkeysalso{#3}}}}
\tikzset{>=latex}
\tikzset{mynodes/.style={circle,white,fill=UniBlue!90,text width=.8em,inner sep=0pt,text centered,font=\footnotesize}}
\tikzset{myedges/.style={thick}}
\tikzset{mypath/.style={ultra thick, draw=#1}, mypath/.default=AlertOrange}
\tikzset{mysubset/.style={draw,#1,dotted,thick,fill=#1!10}, mysubset/.default=UniBlue}
\tikzset{mypath2/.style={ultra thick, dashed, draw=UniGreen}}
\tikzset{graphs/mygraph/.style={nodes={mynodes}, edges={myedges}, empty nodes}}
\tikzset{nomargin/.style={inner sep=0pt, outer sep=0pt}}
%% graph macros
\tikzgraphsset{declare={mybipartite}{%
subgraph I_nm [V={1,2,3,4,5,6}, W={a,b,c,d,e,f}];
1 -- {a,b};
2 -- {b,c,d,e};
3 -- {d,f};
4 -- {e,f};
{5,6} -- f;
}}
\tikzgraphsset{declare={mybipartitedi}{%
subgraph I_nm [V={1,2,3,4,5,6}, W={a,b,c,d,e,f}];
1 -> {a,b};
2 -> {b,c,d,e};
3 -> {d,f};
4 -> {e,f};
{5,6} -> f;
}}
\tikzgraphsset{declare={mybipartite2}{%
subgraph I_nm [V={1,2,3}, W={a,b,c}];
1 -- {a,b}; 2 -- {a,b,c}; 3 -- {a,c};
}} 
% Petersen graph
\tikzgraphsset{declare={Petersen}{%
    subgraph C_n [n=5,name=A, radius=1.5cm]; 
    subgraph I_n [n=5,name=B, radius=.8cm];
    \foreach \i [evaluate={\j=int(mod(\i+2,5)+1);}] in {1,...,5}{
        A \i -- B \i;
        B \i -- B \j;
}}}
% Example graph for Menger's theorem
\tikzgraphsset{declare={menger}{%
    s[at={(-1,0)}, as={$s$}];
    (s) -> 1[at={(0,1)}] -> 2[at={(1,0)}];
    (s) -> 3[at={(0,-1)}] -> 2;
    4[at={(2,1)}] <- 5[at={(2,-1)}];
    {5,6[at={(3,1)}], 7[at={(3,-1)}]} -> t[at={(4,0)}, as={$t$}];
    (s) -> (2);
    (2) -> {(4), (5)};
    {(4),(5)} -> {(6),(7)};
    (6) -> (7);
    1 -> 4;
}} 
\tikzgraphsset{declare={undirected menger}{%
    s[at={(-1,0)}, as={$s$}];
    (s) -- 1[at={(0,1)}] -- 2[at={(1,0)}];
    (s) -- 3[at={(0,-1)}] -- 2;
    4[at={(2,1)}] -- 5[at={(2,-1)}];
    {5,6[at={(3,1)}], 7[at={(3,-1)}]} -- t[at={(4,0)}, as={$t$}];
    (s) -- (2);
    (2) -- {(4), (5)};
    {(4),(5)} -- {(6),(7)};
    (6) -- (7);
    1 -- 4;
}} 

% Appendix
\usepackage{appendixnumberbeamer}

% macros
\newcommand{\R}{\mathbb{R}}
\newcommand{\Z}{\mathbb{Z}}
\newcommand{\be}{\mathbf{e}}
\newcommand{\zeros}{\mathbf{0}}
\newcommand{\ones}{\mathbf{1}}
\newcommand{\bfzero}{\zeros}
\newcommand{\bfone}{\ones}

% \newcommand{\bp}{\mathbf{p}}
\newcommand{\eps}{\varepsilon}
\newcommand{\defiff}{\overset{\text{def}}{\iff}}
\DeclareMathOperator{\poly}{poly}
\DeclareMathOperator{\conv}{conv}
\DeclareMathOperator{\In}{In}
\DeclareMathOperator{\Out}{Out}
\DeclareMathOperator{\val}{val}
\newcommand{\lo}{{\mathrm{lo}}}
\newcommand{\up}{{\mathrm{up}}}
\newcommand{\caI}{\mathcal{I}}
\newcommand{\caB}{\mathcal{B}}
\newcommand{\caC}{\mathcal{C}}
\newcommand{\bM}{\mathbf{M}}
\newcommand{\IO}{\mathrm{IO}}

\usepackage{mathtools}
\DeclarePairedDelimiter{\abs}{\lvert}{\rvert}
\DeclarePairedDelimiter{\norm}{\lVert}{\rVert}
\DeclarePairedDelimiter{\inprod}{\langle}{\rangle}
\usepackage[no-test-for-array]{nicematrix}
\usepackage{diagbox}

\newcommand{\highlightcap}[3][red]{\tikz[baseline=(x.base)]{\node[rectangle,rounded corners,fill=#1!10](x){#2} node[below=0pt of x, color=#1]{#3};}}
\newcommand{\highlight}[2][red]{\tikz[baseline=(x.base)]{\node[rectangle,rounded corners,fill=#1!10](x){#2};}}
\newcommand{\mycite}[2][\footnotesize]{\textcolor{UniBlue}{#1 [#2]}}

%section
% \AtBeginSection[]{
%     \frame[c]{
%         \centering
%         \begin{tikzpicture}
%             \node[white, circle, fill=UniBlue, text centered, text width=.3\linewidth](bullet) {
%                 {\Huge \bf \insertsectionnumber} \\[4em]
%     };
%     \node[text width=.65\linewidth, right=1em of bullet]{%
%         \usebeamerfont{title}\usebeamercolor[fg]{title}\insertsection\par %
%     };
%         \end{tikzpicture}
%     } %目次スライド
% }

%grid
% \usepackage[colorgrid,gridunit=pt,texcoord]{eso-pic}
\newcommand{\postit}[1]{\tikz[overlay, remember picture]{\node[draw,anchor=north east, align=center, rectangle, fill=yellow!10, xshift=-1pt, yshift=-1pt] at (current page.north east) {#1};}}

\usepackage{tcolorbox}
\newtcolorbox{primal}{colback=red!10!white, colframe=red, title=主問題 (P), left=2pt,right=2pt,top=0pt,bottom=0pt}
\newtcolorbox{dual}{colback=blue!10!white, colframe=blue, title=双対問題 (D), left=2pt,right=2pt,top=0pt,bottom=0pt}

% visible on
\usetikzlibrary{overlay-beamer-styles}
\usepackage[normalem]{ulem}
\usepackage{pxrubrica}
\usepackage{booktabs}

\newcommand{\Aphiup}{\textcolor{UniBlue}{A_\varphi^\up}}
\newcommand{\Aphilow}{\textcolor{AlertOrange}{A_\varphi^\lo}}
\newcommand{\bv}{\boldsymbol{v}}
\newcommand{\bx}{\boldsymbol{x}}
\newcommand{\by}{\boldsymbol{y}}
\newcommand{\bz}{\boldsymbol{z}}
\newcommand{\bw}{\boldsymbol{w}}
\newcommand{\bc}{{\boldsymbol c}}
\DeclareMathOperator{\supp}{supp}
\DeclareMathOperator{\argmax}{argmax}
\DeclareMathOperator{\argmin}{argmin}
\DeclareMathOperator{\E}{\mathbb{E}}
\csname endofdump\endcsname
\title{組合せ最適化特論}
\subtitle{第7回（マトロイド）}
\author{担当: \textbf{相馬 輔}}
\date{2026/7/2}

\begin{document}
\maketitle

\begin{frame}{}
    \centering\small
    \tikzset{block/.style={rectangle, draw=gray, fill=gray!10, thick, minimum width=3em, minimum height=2em, align=center}}
    \begin{tikzpicture}
        \node[block,] (BM) {重みなし\\二部マッチング};
        \node[block,above=1em of BM.north west, xshift=-2em] (WBM) {重みつき\\二部マッチング};
        \node[block,above=5em of BM.north east, xshift=1em] (NF) {最大流};
        \node[block,above=9em of BM] (MCF) {最小費用流};
        % \node[block,above=1em of NF.north east, xshift=1em] (MC) {最小カット};
        \node[block,right=11em of BM] (MTR) {マトロイド};
        \node[block,above=7em of MTR] (MI) {マトロイド交差};
        \node[block,above=3em of MCF, xshift=7em] (SFM) {劣モジュラ最小化};
        \graph[mygraph]{
            (BM) -> {(WBM), (NF)} -> (MCF);
            {(MTR), (BM)} -> (MI);
            % (NF) -> (MC);
            {(MCF), (MI)} -> (SFM);
        };
        \begin{scope}[on background layer]
            \node[mysubset, fit=(BM) (WBM), "二部マッチング" {below, UniBlue}]{};
            \node[mysubset, AlertOrange, fill=AlertOrange!10, fit=(NF) (MCF), "ネットワークフロー" {above left, AlertOrange}]{};
            \node[mysubset, UniGreen, fill=UniGreen!10, fit=(MTR) (MI), "マトロイド" {below, UniGreen}]{};
        \end{scope}
    \end{tikzpicture}
\end{frame}

\begin{frame}[allowframebreaks]
    \frametitle{各回の内容（予定）} 
    \small
    \begin{enumerate}
        \item (4/23) イントロ+多面体的組合せ論（線形計画法の復習，整数多面体，完全単模行列）
        \item[] \textcolor{gray}{(4/30) 休み}
        \item (5/7) 二部マッチング①（Konig-Egervaryの定理，増加道アルゴリズム，ハンガリー法）
        \item (5/14) 二部マッチング②（最短路問題の復習，逐次最短路法と主双対法，最適性基準からの見方）
        \item (5/21) 最大流① （定式化，最大流最小カット定理，応用例）
        \item (5/28) 最大流②（残余ネットワーク，Ford--Fulkerson法，Edmonds--Karpのアルゴリズム）+ 最小費用流①（定式化，輸送問題，最大流との関係）
        \item[] \textcolor{gray}{(6/4) 休み}
        % \item[] \sout{(12/11) 最小カット（永持茨木のアルゴリズム，Kargerのアルゴリズム）}
        \item (6/11) 最小費用流②（逐次最短路法，容量スケーリング法）
        \item[] \textcolor{gray}{(6/18) 休み}
        \item[] \textcolor{gray}{(6/25) 休み}
        \item (7/2) マトロイド（定義と公理系，貪欲法，マトロイド多面体） 
        \item (7/9) 独立マッチング・マトロイド交差 （定義，応用例，Edmondsの最大最小定理，増加道アルゴリズム）
        \item (7/16) 劣モジュラ関数①（諸例，劣モジュラ基多面体，Lov\'asz拡張，劣モジュラ最小化）
        \item[] \textcolor{gray}{(7/23) 休み}
        \item (7/30) 劣モジュラ関数②（劣モジュラ最大化，近似アルゴリズム，貪欲法）
        \item (どこか) 精選トピック①
        \item (どこか) 精選トピック②
        \item (どこか) 精選トピック③
    \end{enumerate}

    \bigskip
    \begin{itemize}
        \item レポート課題を6月と期末に出す予定
        \item 精選トピックの日時は今後調整 
    \end{itemize}
\end{frame}

\begin{frame}<1>[label=outline]{目次}
\begin{tikzpicture}
    \tikzstyle{block} = [rectangle, text width=\textwidth, rounded corners, minimum height=4em]
    \tikzstyle{line} = [draw, -latex]

    \node[block,alt={<1,3->{fill=cyan!10}{fill=red!10}}] (combopt) {
        \begin{columns}
            \begin{column}{0.5\linewidth}
        \textbf{1. マトロイドとは}
            \end{column}
            \begin{column}{0.4\linewidth}
                \small
                \structure{・} マトロイドの諸例 \\
                \structure{・} 独立集合族 \\
                \structure{・} 基族 \\
                \structure{・} 階数関数 \\
            \end{column}
        \end{columns}
    };
    \node[block,alt={<1-2,4->{fill=cyan!10}{fill=red!10}}, below=.5em of combopt] (LP) {
        \begin{columns}
            \begin{column}{0.5\linewidth}
        \textbf{2. 最小重み独立集合と貪欲法}
            \end{column}
            \begin{column}{0.4\linewidth}
                \small
                \structure{・} 最小重み独立集合問題 \\
                \structure{・} 貪欲法 \\
                \structure{・} Kruskalのアルゴリズム
            \end{column}
        \end{columns}
    };
    \node[block,alt={<1-3,5->{fill=cyan!10}{fill=red!10}}, below=.5em of LP] {
        \begin{columns}
            \begin{column}{0.5\linewidth}
        \textbf{3. マトロイド多面体}
            \end{column}
            \begin{column}{0.4\linewidth}
                \small
                \structure{・} 線形不等式系と整数性 \\
            \end{column}
        \end{columns}
    };
\end{tikzpicture}%
\end{frame}


%%%%%%%%%%%%%%%%%%%%%%%%%%%%%%%%%%%%%%%%%%%%%
\section{マトロイドとは}
\againframe<2|handout:0>{outline}
%%%%%%%%%%%%%%%%%%%%%%%%%%%%%%%%%%%%%%%%%%%%%

\begin{frame}{マトロイドとは}
    \small

    % {\it \alert{マトロイド (matroid)}  (matrix~+~-oid)} \\
    \begin{block}{マトロイド (matroid; matrix~+~-oid)}
    Whitney (1935), 中澤(1935), Pauc--Haupt--N\"obeling (1937--1940), Rado (1942) らによって独立に導入された組合せ構造．
    \end{block}

    \vfill
    \structure{組合せ最適化におけるマトロイドの重要性}\\
    基本的な組合せ最適化問題にマトロイド制約を加えることで，非自明な一般化が得られる．
    \begin{itemize}
        \item \textbf{二部マッチング $\longrightarrow$ 独立マッチング・マトロイド交差（次回）}
        \item サイズ制約劣モジュラ最大化 $\longrightarrow$ マトロイド制約劣モジュラ最大化
        \item 秘書問題 $\longrightarrow$ マトロイド秘書問題
        \item[] ...
    \end{itemize}
\end{frame}

\newcommand{\indepaxiominner}{%
        \item[(I0)] $\emptyset \in \caI$
        \item[(I1)] $I \in \caI$, $I' \subseteq I$ $\implies$ $I' \in \caI$
        \item[(I2)] $I_1, I_2 \in \caI$, $\abs{I_1} < \abs{I_2}$ $\implies$ $\exists e \in I_2 \setminus I_1$~s.t.~$I_1 + e \in \caI$
}
\newcommand{\indepaxiom}{%
    \begin{description}
        \indepaxiominner
    \end{description}}
\begin{frame}{マトロイドの例①（分割マトロイド）}
    \begin{columns}[T]
    \begin{column}{.8\textwidth}
    $E$: 有限集合 \\
    $(E_1, E_2, E_3)$: $E$の分割 {\footnotesize ※一般に$k$分割でもよい}
    \[
        \caI = \{I \subseteq E: \abs{I \cap E_i} \leq 1 \quad (i=1,2,3) \}
    \]
    \end{column}
    \begin{column}{.2\textwidth}
    \begin{tikzpicture}[yscale=.8]
                \foreach \y in {0,1,2}{
                    \node[circle,fill=UniBlue,inner sep=0pt,text width=.8em] at (0,\y) {};
                    \node[circle,fill=UniPink,inner sep=0pt,text width=.8em] at (1,\y) {};
                    \node[circle,fill=UniGreen,inner sep=0pt,text width=.8em] at (2,\y) {};
                }
                \node[] at (0, 3) {\small $E_1$};
                \node[] at (1, 3) {\small $E_2$};
                \node[] at (2, 3) {\small $E_3$};
                \draw (-.5, -.5) -- (2.5, -.5) -- (2.5, 2.5) -- (-.5, 2.5) -- cycle;
                \draw[dotted, thick] (0.5, -.5) -- (.5, 2.5);
                \draw[dotted, thick] (1.5, -.5) -- (1.5, 2.5);
            \end{tikzpicture}
    \end{column}
    \end{columns}
    
    \vfill
    集合族$\caI$は以下の3つの性質を満たす:
    \begin{block}{}
    \indepaxiom
    \end{block}
    {\footnotesize \hfill ($I_1 + e := I_1 \cup \{e \}$)}
\end{frame}

\begin{frame}{マトロイドの例②（グラフ的マトロイド）}
    \begin{columns}[T]
    \begin{column}{.7\textwidth}
    $G = (V, E)$: 無向グラフ
    \[
        \caI = \{I \subseteq E: \text{$I$は森({\footnotesize 閉路を含まない部分グラフ})}\}
    \]
    \end{column}
    \begin{column}{.3\textwidth}
        \footnotesize
        \centering
        \tikz{\graph[mygraph,clockwise]{
        Petersen;
        {[edges={mypath}] 
            (A 1) -- (A 2) -- (A 3) -- (A 4) -- (A 5);
            (A 3) -- (B 3);
            (B 4) -- (B 2) -- (B 5);
        }
        };
        } \\ 森の例
    \end{column}
    \end{columns}
    \begin{lemma}
    集合族$\caI$は以下の3つの性質を満たす:
    \indepaxiom
    \end{lemma}
\end{frame}

\newcommand{\indepaxiompostit}{%
    \postit{\footnotesize
    \begin{minipage}{0.4\textwidth}
        \begin{description}
        \setlength{\itemindent}{-40pt}
        \setlength{\itemsep}{0pt}
        \indepaxiominner
        \end{description}
    \end{minipage}}
}

% \begin{frame}{補題の証明}
%     \indepaxiompostit 

%     \structure{(証明)} 
%     (I0), (I1)は自明．\\ 
%     (I2): 一般に，$n$頂点の木は$n-1$本の枝を持つことに注意．

%     森$I_1$, $I_2$ ($\abs{I_1} < \abs{I_2}$)を取る．もし，$I_2$の頂点のうち，$I_1$に含まれない頂点$v$が存在すれば，$v$に接続する$I_2$の枝を$e$とすると，$I_1 + e$は森である．
%     \qed
% \end{frame}

\begin{frame}{マトロイドの例③（線形マトロイド）}
    % $v_1, \dots, v_n \subseteq \mathbb{F}^d$: ベクトル空間の元\\
    \begin{columns}[T]
    \begin{column}{.7\textwidth}
    $A$: {\footnotesize （体$\mathbb{F}$上の）}$r \times n$行列 \\ 
    $E = \{1, \dots, n\}$
    \[
        \caI = \{I \subseteq E: \text{列部分行列$A[I]$の列は一次独立}\}
    \]
    \end{column}
    \begin{column}{.3\textwidth}
        \centering
        \begin{tikzpicture}
            \coordinate (O) at (0,0);
            \coordinate (A) at (4,0);
            \coordinate (B) at (0,2);
            \coordinate (Il) at ($(B)!0.2!(A |- B)$);
            \coordinate (Ir) at ($(B)!0.6!(A |- B)$);
            \begin{scope}[on background layer]
                \path[fill=AlertOrange!10] (Il) rectangle (Ir |- O);
                \draw[rectangle, thick] (O) rectangle (A |- B);
            \end{scope}
            \draw[mypath, <->, AlertOrange] (Il) -- (Ir) node[midway, above] {\small $I$};
            \node[AlertOrange] at ($(Il)!0.5!(Ir |- A)$) {\small $A[I]$};
        \end{tikzpicture}
    \end{column}
    \end{columns}

    \begin{lemma}
    集合族$\caI$は以下の3つの性質を満たす:
    \indepaxiom
    \end{lemma}
\end{frame}

\begin{frame}{補題の証明}
    \indepaxiompostit 

    \begin{proof}
    \textbf{(I0), (I1):} 定義から自明． 

    \bigskip
    \textbf{(I2):} $I_1, I_2 \in \caI$ ($\abs{I_1} < \abs{I_2}$)を取る．
    
    \bigskip
    \structure{Claim.} $A[I_2]$の列のうち，$A[I_1]$の列が張る部分空間に含まれないものが存在する．\\
    {\small \structure{(∵)} $A[I_1]$, $A[I_2]$の列の張る部分空間をそれぞれ$W_1, W_2$とすると，$\dim W_1 = \abs{I_1} < \abs{I_2} = \dim W_2$であるから．\qed}
    
    \bigskip
    そのような列の添字を$e \in I_2 \setminus I_1$とすると，$I_1 + e \in \caI$.
    \end{proof}
\end{frame}

\begin{frame}{分割マトロイド$\subset$線形マトロイド}

    \begin{columns}[T]
    \begin{column}{.8\textwidth}
    $E$: 有限集合 \\
    $(E_1, E_2, E_3)$: $E$の分割 {\footnotesize ※一般に$k$分割でもよい} \\
    に対して，行列
    \begin{align*}
        A = \begin{bNiceArray}{ccc|ccc|ccc}[first-row]
        \Block{1-3}{E_1} &&& \Block{1-3}{E_2} &&& \Block{1-3}{E_3} \\
        1 & 1 & 1 & 0 & 0 & 0 & 0 & 0 & 0 \\
        0 & 0 & 0 & 1 & 1 & 1 & 0 & 0 & 0 \\
        0 & 0 & 0 & 0 & 0 & 0 & 1 & 1 & 1 \\
        \end{bNiceArray}
    \end{align*}
    を考えると，$A$が定める集合族$\caI$は分割マトロイドの集合族$\caI$に一致する．
    \end{column}
    \begin{column}{.2\textwidth}
    \begin{tikzpicture}[yscale=.8]
                \foreach \y in {0,1,2}{
                    \node[circle,fill=UniBlue,inner sep=0pt,text width=.8em] at (0,\y) {};
                    \node[circle,fill=UniPink,inner sep=0pt,text width=.8em] at (1,\y) {};
                    \node[circle,fill=UniGreen,inner sep=0pt,text width=.8em] at (2,\y) {};
                }
                \node[] at (0, 3) {\small $E_1$};
                \node[] at (1, 3) {\small $E_2$};
                \node[] at (2, 3) {\small $E_3$};
                \draw (-.5, -.5) -- (2.5, -.5) -- (2.5, 2.5) -- (-.5, 2.5) -- cycle;
                \draw[dotted, thick] (0.5, -.5) -- (.5, 2.5);
                \draw[dotted, thick] (1.5, -.5) -- (1.5, 2.5);
            \end{tikzpicture}
    \end{column}
    \end{columns}
\end{frame}

\begin{frame}{グラフ的マトロイド$\subset$線形マトロイド}
    \small

    \bigskip
    \begin{columns}
    \begin{column}{.7\textwidth}
    \structure{ネットワーク行列 (復習)} 有向グラフ$G = (V, A)$に対し，\\ 
    $V \times A$行列$B$を
    \[
        B_{i,a} := \begin{cases}
            +1 & \text{($i$は$a$の始点)} \\
            -1 & \text{($i$は$a$の終点)} \\
            0  & \text{otherwise}
        \end{cases} 
    \]
    $(i \in V, a \in A)$と定める．
    \end{column}
    \begin{column}{.3\textwidth}
        \centering
        \begin{tikzpicture}
            \graph[mygraph, math nodes, grow right=1.4cm, branch down=1.4cm]{
            a -> b;
            c -> d;
            a -> {c, d};
            b -> d;
        };
        \end{tikzpicture}
        \footnotesize
        \[
            B = \begin{bmatrix}
            +1 & +1 & +1 & 0 & 0 \\
            -1 & 0 & 0 & +1 & 0 \\
            0 & -1 & 0 & 0 & +1 \\
            0 & 0 & -1 & -1 & -1
            \end{bmatrix}
        \]
    \end{column}
    \end{columns}

    \vfill
    \begin{lemma}
        無向グラフ$G = (V, E)$の枝を任意に向き付けた有向グラフを考え，そのネットワーク行列を$B$とする．
        このとき，
        $B[I]$の列は一次独立 $\iff$ $I$は$G$の森
        % 無向グラフ$G$が定めるグラフ的マトロイドと，ネットワーク行列$B$が定める線形マトロイドは一致する．
    \end{lemma}
\end{frame}

\begin{frame}{マトロイドの定義（独立集合族）}
    \begin{definition}
        $E$: 有限集合, $\caI \subseteq 2^E$: 部分集合族 \\
        $\bM = (E, \caI)$が\alert{マトロイド} $\defiff$
        \indepaxiom
    
        \bigskip
        特に，$I \in \caI$を\alert{独立集合 (independent set)}といい，$\caI$を\alert{独立集合族}という．
        また，$E$をマトロイドの\alert{台集合 (ground set)}という．
    \end{definition}


    \tikz[overlay, remember picture]{
        \node[nomargin] at ([xshift=-2.5cm, yshift=1.1cm]current page.east) {
        \tikz{
            \node[circle, fill=UniBlue, opacity=.3, minimum size=1.5cm] (B1) at (0,0) {};
            \node[circle, fill=AlertOrange, opacity=.3, minimum size=2cm] (B2) at (1,0) {};
            \node[point,fill=AlertOrange,"$e$" above] (j) at ([xshift=+.4cm] B2.base) {};
            \node[UniBlue, above=1pt of B1, xshift=-.4cm] {$I_1$};
            \node[AlertOrange, above=1pt of B2, xshift=+.4cm] {$I_2$};
        }
            };
    }
\end{frame}

\begin{frame}{基族}
    $\bM = (E, \caI)$: マトロイド
    \begin{definition}
    \begin{itemize}
        \item 独立集合$I$が\alert{極大 (maximal)} $\defiff$ $I$を真に含む独立集合が存在しない
        \item  極大な独立集合を$\bM$の\alert{基 (base)}といい，基の全体を\alert{基族 (base family)}という．        
    \end{itemize}
    \end{definition}

    \bigskip
    \small
    \structure{例}
    \begin{itemize}
        \item 分割マトロイドの基族: 各分割ブロック$E_i$から要素を1つずつ含む集合
        \item グラフ的マトロイドの基族: グラフが連結であるとすると，基はグラフの\alert{全域木 (spanning tree)}
        % \item 線形マトロイドの基族: 
    \end{itemize}
\end{frame}
    
\begin{frame}{基族}
    \bigskip
    $\caB$をマトロイドの基族とする．
    \begin{lemma}
        $B_1, B_2 \in \caB \implies \abs{B_1} = \abs{B_2}$
    \end{lemma}
    \pause
    {\small 
    \begin{proof}
    背理法．$\abs{B_1} < \abs{B_2}$とする．
    $B_1, B_2$は独立集合なので，公理(I2)より，ある$e \in B_2 \setminus B_1$で$B_1 + e \in \caI$となるものが存在．これは$B_1$の極大性に反する．
    \end{proof}
    }

    \pause
    \bigskip
    基の大きさをマトロイド$\bM$の\alert{階数 (rank)}といい，$r(\bM)$で表す．
    \begin{lemma}
        $I \in \caI$, $\abs{I} = r(\bM)$ $\implies$ $I$は基．
    \end{lemma}
\end{frame}

\begin{frame}{基族}
    \begin{columns}[T]
    \begin{column}{.75\textwidth}
    \begin{theorem}
        マトロイドの基族$\caB$に対し，次が成り立つ:
    \begin{description}
        \item[(B1)] $\caB \neq \emptyset$
        \item[(B2)] $B_1, B_2 \in \caB$, $B_1 \neq B_2$, $i \in B_1 \setminus B_2$ \\ 
        $\implies \exists j \in B_2 \setminus B_1$~s.t.~$B_1 - i + j \in \caB$.
    \end{description}
    \end{theorem}
    \end{column}
    \begin{column}{.25\textwidth}
        \centering
           \vspace{1em}
           \begin{tikzpicture}
            %    \path[fill=UniBlue, opacity=.3] (0,0) circle[radius=1];
            %    \path[fill=AlertOrange, opacity=.3] (1,0) circle[radius=1];
               \node[circle, fill=UniBlue, opacity=.3, minimum size=2cm] (B1) at (0,0) {};
               \node[circle, fill=AlertOrange, opacity=.3, minimum size=2cm] (B2) at (1,0) {};
               \node[point,fill=UniBlue,"$i$" left] (i) at ([xshift=-.4cm] B1.base) {};
               \node[point,fill=AlertOrange,"$j$" right] (j) at ([xshift=+.4cm] B2.base) {};
               \node[UniBlue, above=1pt of B1, xshift=-.4cm] {$B_1$};
               \node[AlertOrange, above=1pt of B2, xshift=+.4cm] {$B_2$};
           \end{tikzpicture} 
    \end{column}
    \end{columns}

    \vfill
    \small
    \pause
    \begin{proof} \\
    \textbf{(B1):} 公理(I0) $\emptyset \in \caI$より，少なくとも1つは基が存在する．

    \smallskip
    \textbf{(B2):} $B_1 - i$と$B_2$は独立集合で，$\abs{B_1 - i} = \abs{B_1} - 1 < \abs{B_2}$．よって，公理(I2)より，ある$j \in B_2 \setminus (B_1 - i) = B_2 \setminus B_1$で$(B_1 - i) + j \in \caI$となるものが存在．$\abs{B_1 - i + j} = \abs{B_1}$なので，補題よりこれは基．
    \end{proof}
\end{frame}

\begin{frame}{基族}
    実は，基族の性質(B1), (B2)からマトロイドを定義することもできる．

    \begin{lemma}[cf.~\mycite{Ch.~14, Korte--Vygen (2018)}]
        $\caB \subseteq 2^E$を(B1), (B2)を満たす集合族とする．このとき，
        \[ 
            \caI := \{ I \subseteq E: \text{$I$はある$B \in \caB$に含まれる} \}
        \]
        とすると，$\caI$は独立集合族である．
    \end{lemma}

    すなわち，(B1), (B2)もマトロイドの公理系である．
\end{frame}

\begin{frame}{階数関数}
    $\bM = (E, \caI)$: マトロイド
    \begin{definition}
        \begin{itemize}
            \item $X \subseteq E$に対し，$X$に含まれる最大の独立集合の大きさを$X$の\alert{階数 (rank)}といい，$r(X)$で表す．
            \item $r: X \mapsto r(X)$を\alert{階数関数 (rank function)}という．
        \end{itemize}
        % $\bM$の\alert{階数関数 (rank function)} $r: 2^V \to \Z_+$を次のように定める:
        % \[
        %     r(X) := \max\{\abs{B \cap X}: B \in \caB \}
        %     \qquad \text{\small $X$に含まれる最大の独立集合の大きさ}
        % \]
    \end{definition}

    \pause
    \bigskip
    \structure{例}
    \begin{itemize}
        \item 分割マトロイド: $r(X) = \abs{\{i: E_i \cap X \neq \emptyset\}}$
        \item グラフ的マトロイド: $r(X) = \abs{V(G[X])} - \text{($G[X]$の連結成分数)}$ \\ 
              {\footnotesize \hfill ($G[X]$: $X$の誘導部分グラフ)}
        \item 線形マトロイド: $r(X) = \rank A[X]$
    \end{itemize}
\end{frame}

\begin{frame}{階数関数の性質}
    \begin{theorem}
        マトロイドの階数関数$r$は次の性質を満たす:
        \begin{description}
            \item[(R1)] $0 \leq r(X) \leq \abs{X}$ \quad ($X \subseteq E$) 
            \item[(R2)] $r(X) \leq r(Y)$ \quad ($X \subseteq Y \subseteq E$) \hfill \alert{単調性}
            \item[(R3)] $r(X) + r(Y) \geq r(X \cup Y) + r(X \cap Y)$ \quad ($X, Y \subseteq E$) \\ \hfill\alert{劣モジュラ性}
        \end{description}
    \end{theorem}

    \vfill
    \small
    実は，(R1)--(R3)もマトロイドの公理系となる．
    \begin{theorem}[cf.~\mycite{Ch.~14, Korte--Vygen (2018)}]
        (R1)--(R3)を満たす集合関数$r: 2^E \to \Z_+$に対し，
        \setlength{\abovedisplayskip}{1pt}
        \setlength{\belowdisplayskip}{1pt}
        \[
            \caI = \{X \subseteq E: \abs{X} = r(X) \} 
        \]
        はマトロイドの独立集合族．
    \end{theorem}
\end{frame}

\begin{frame}{定理の証明}
    \postit{\begin{minipage}{.6\textwidth}\footnotesize 
        \begin{description}
        \setlength{\itemindent}{-35pt}
        \setlength{\leftmargin}{0pt}
        \setlength{\itemsep}{0pt}
            \item[(R1)] $0 \leq r(X) \leq \abs{X}$
            \item[(R2)] $r(X) \leq r(Y)$ \quad ($X \subseteq Y \subseteq E$) 
            \item[(R3)] $r(X) + r(Y) \geq r(X \cup Y) + r(X \cap Y)$ \quad ($X, Y \subseteq E$) 
        \end{description}
    \end{minipage}}

    \small
    \begin{proof}
    \textbf{(R1), (R2):} 定義から自明．  

    \begin{columns}[T]
    \begin{column}{.8\textwidth}
    \textbf{(R3):} $I_{X\cap Y}$を$X\cap Y$の最大独立集合とする．
    公理(I2)より，$I_{X \cap Y}$に$X \cup Y$の要素を加えて，$X \cup Y$の最大独立集合$I_{X \cup Y}$を作れる．
    
    % \begin{align*}
    %     I_{X\cap Y} = I_{X\cup Y} \cap (X \cap Y)
    %     % , \quad
    %     % I_{X\cup Y} = B^* \cap (X \cup Y)
    % \end{align*}
    このとき，
    \begin{itemize}
        \item $I_{X\cap Y} = I_{X\cup Y} \cap (X \cap Y)$
        \item $I_{X \cup Y} \cap X$は$X$の独立集合，$I_{X \cup Y} \cap Y$は$Y$の独立集合．
    \end{itemize}
    \end{column}
    \begin{column}{.2\textwidth}
        \centering
           \begin{tikzpicture}
               \node[circle, fill=UniBlue, opacity=.3, minimum size=2cm] (B1) at (0,0) {};
               \node[circle, fill=AlertOrange, opacity=.3, minimum size=2cm] (B2) at (1,0) {};
               \node[UniBlue, above=1pt of B1, xshift=-.4cm] {$X$};
               \node[AlertOrange, above=1pt of B2, xshift=+.4cm] {$Y$};
               \node at ($(B1)!0.5!(B2)$) {\footnotesize $X \cap Y$};
           \end{tikzpicture} 
    \end{column}
    \end{columns}

    \bigskip
    よって，
    \begin{align*}
        r(X) + r(Y) 
        &\geq |I_{X \cup Y} \cap X| + |I_{X \cup Y} \cap Y| \\
        &= |I_{X \cup Y} \cap (X \cup Y) | + |I_{X \cup Y} \cap (X \cap Y)| \\
        &= |I_{X \cup Y}| + |I_{X \cap Y}| \\
        &= r(X \cup Y) + r(X \cap Y) 
    \end{align*}
    \end{proof}
\end{frame}

%%%%%%%%%%%%%%%%%%%%%%%%%%%%%%%%%%%%%%%%%%%%%
\section{最小重み独立集合と貪欲法}
\againframe<3|handout:0>{outline}
%%%%%%%%%%%%%%%%%%%%%%%%%%%%%%%%%%%%%%%%%%%%%
\begin{frame}{最小重み独立集合問題}
    \begin{itemize}
        \item $E$上のマトロイド$\bM = (E, \caI)$ %\quad ($\caB$: $\bM$の基族) 
        \item 台集合上の重み$w: E \to \R$
        \item 非負整数$0 \leq k \leq r(\bM)$
    \end{itemize}
    
    \vfill
    \begin{block}{最小重み独立集合問題 (Minimum Weight Independent Set Problem)}
        \[
           \text{minimize} \quad w(I) \quad\text{s.t.}\quad I \in \caI, \abs{I} = k
        \]
    \end{block}

    \vfill
    \structure{この節の目標:} 最小重み独立集合問題は\alert{貪欲法 (greedy algorithm)}で解けることを示す．
\end{frame}

\begin{frame}{例：最小重み全域木問題}

    \begin{columns}[T]
        \begin{column}{.8\textwidth}
            \begin{itemize}
                \item $G = (V, E)$ 連結無向グラフ
                \item $w: E \to \R$ 枝重み
            \end{itemize}
    \begin{block}{最小重み全域木問題 (Minimum Weight Spanning Tree Problem)}
        \[
           \text{minimize} \quad w(T) \quad\text{s.t.}\quad \text{$T$は$G$の全域木} 
        \]
    \end{block} 
        \end{column}
    \begin{column}{.2\textwidth}
        \centering
        \tikz{\graph[mygraph,clockwise]{
        Petersen;
        {[edges={mypath}] 
            (A 1) -- (A 2) -- (A 3) -- (A 4) -- (A 5);
            (A 3) -- (B 3) -- {(B 1), (B 5)} ;
            (A 4) -- (B 4) -- (B 2);
        }
        };
        } \\ 全域木の例
    \end{column}
    \end{columns}

    これは
    \begin{itemize}
        \item $\bM =$ $G$のグラフ的マトロイド
        \item $k = |V| - 1 = r(\bM)$
    \end{itemize}
    に対する最小重み独立集合問題．
\end{frame}

\begin{frame}<1,8>[label=MST]
    \frametitle{例}

    \begin{center}
        \tikz[scale=2]{
            \graph[mygraph, no placement]{
            a[at={(0, 0)}];
            b[at={(1, 1)}];
            c[at={(3, 1)}];
            d[at={(5, 1)}];
            e[at={(6, 0)}];
            f[at={(5,-1)}];
            g[at={(3,-1)}];
            h[at={(1,-1)}];
            i[at={(2, 0)}];
            };
            \foreach \i/\j/\w/\pos in {a/b/4/above, b/c/8/above, c/d/7/above, d/e/9/above, e/f/10/below, f/g/2/below, g/h/1/below, h/a/8/below, b/h/11/left, h/i/7/above, i/c/2/below, i/g/6/below, c/f/4/below, d/f/14/right}{
                  \draw[myedges] (\i) edge ["{$\w$}" \pos] (\j);
            }

            \begin{scope}[every path/.style={mypath}]
                \path<2-> (h) -- (g);
                \path<3-> (g) -- (f);
                \path<4-> (i) -- (c);
                \path<5-> (a) -- (b);
                \path<6-> (c) -- (f);
                \path<7-> (c) -- (d);
                \path<8-> (c) -- (d) -- (e) (a) -- (h);
            \end{scope}

        }
    \end{center}
\end{frame}

\begin{frame}{貪欲法}

    \begin{block}{}
    \begin{algorithmic}[1]
        \STATE $X_0 := \emptyset$, $k := 0$.
        \STATE $E$の要素を$w$の昇順に並べる: $w(e_1) \leq \dots \leq w(e_n)$ ($n = \abs{E}$).
        \FOR{$i = 1, \dots, n$}
            \STATE $X_k + e_i \in \caI$ならば$e_i$を$X_k$に追加し，$k \gets k + 1$とする． \\
            \COMMENT{$X_k$は大きさ$k$の独立集合}
        \ENDFOR
    \end{algorithmic}
    \end{block}

    \vfill
    \begin{theorem}
        $k = 0, 1, \dots, r(\bM)$に対し，$X_k$が存在し，かつ$X_k$は大きさ$k$の独立集合のうち重み最小である．
    \end{theorem}
\end{frame}

\begin{frame}{Kruskalのアルゴリズム}
    特に最小重み全域木問題に適用したものは\alert{Kruskalのアルゴリズム}と呼ばれる．
    
    \begin{block}{}
    \begin{algorithmic}[1]
        \STATE $X := \emptyset$.
        \STATE 枝を$w$の昇順に並べる: $w(e_1) \leq \dots \leq w(e_m)$ ($m = \abs{E}$).
        \FOR{$i = 1, \dots, m$}
            \STATE $X + e_i$が閉路を含まないならば$e_i$を$X$に追加する．
        \ENDFOR
        \RETURN $X$
    \end{algorithmic}
    \end{block}

    \vfill
    \begin{corollary}
        $X$は最小重み全域木である．
    \end{corollary}
\end{frame}

\againframe<1-8>{MST}

\begin{frame}{定理の証明}
    \small
    \begin{lemma}
        貪欲法の任意の時点の暫定解$X$に対して，次が成り立つ:\\
        $e \notin X$, $X + e \in \caI \implies$ $e$はまだ調べていない要素．
    \end{lemma}
    \pause
    \begin{proof}
    $e$は既に調べた要素だと仮定する．
    $e$を調べたときの暫定解を$X' \subseteq X$とすると，$X' + e \subseteq X + e \in \caI$より，$X' + e \in \caI$.
    よって，貪欲法は$e$を$X'$に追加したことになり，$e \in X$となって矛盾．
    \end{proof}

    \vfill
    \pause
    \begin{lemma}
        $k = 0, 1, \dots, r(\bM)$に対し，$X_k$が存在する．
    \end{lemma}
    \pause
    \begin{proof}
    $k < r(\bM)$とし，$X_k$が存在すると仮定する．
    % $B$を$X_k$を含む基とする．
    $k = \abs{X_k} < r(\bM)$より，$X_k$は基ではないから，$e \notin X_k$で$X_k + e \in \caI$となるものが存在．

    \pause
    前補題より，$e$は$X_k$が定義された時点で，まだ調べていない要素である．
    よって，少なくとも1つは$X_k$に追加できる要素がまだ残っているので，$X_{k+1}$が定義される． 
    \end{proof}
\end{frame}

\begin{frame}{定理の証明}
    \small
    \begin{lemma}
        $k = 0, 1, \dots, r(\bM)$に対し，$X_k$は大きさ$k$の最小重み独立集合．
    \end{lemma}

    \pause
    \structure{(証明)}
    背理法．$X_k$より真に重みの小さい大きさ$k$の独立集合$Y$が存在すると仮定．
    \begin{alignat*}{2}
    X_k &= \{x_1, \dots, x_k\} &\quad& (w(x_1) \leq \dots \leq w(x_k)) \\
    Y   &= \{y_1, \dots, y_k\} &\quad& (w(y_1) \leq \dots \leq w(y_k)) 
    \end{alignat*}
    とする．

    \bigskip
    $w(Y) < w(X_k)$より，ある$1 \leq i \leq k$で$w(y_i) < w(x_i)$が成り立つ．
    % そのような$i$のうち最小のものを取る．
    
    \begin{center}
        \begin{tikzpicture}
            \matrix (m) [matrix of nodes, ampersand replacement=\&]{
                $x_1$ \& $\leq\dots\leq$ \& $x_{i-1}$ \& $\leq$ \& $x_i$ \& $\leq\dots\leq$ \& $x_k$ \\[1em]
                $y_1$ \& $\leq\dots\leq$ \& $y_{i-1}$ \& $\leq$ \& $y_i$ \& $\leq\dots\leq$ \& $y_k$ \\
            };
            \node[AlertOrange, rotate=90] at ($(m-1-5)!0.5!(m-2-5)$) {$<$};
        \end{tikzpicture}
    \end{center}
\end{frame}


\begin{frame}{定理の証明（続き）}
    \small
    \begin{center}
        \begin{tikzpicture}
            \matrix (m) [matrix of nodes, ampersand replacement=\&]{
                $x_1$ \& $\leq\dots\leq$ \& $x_{i-1}$ \& $\leq$ \& $x_i$ \& $\leq\dots\leq$ \& $x_k$ \\[1em]
                $y_1$ \& $\leq\dots\leq$ \& $y_{i-1}$ \& $\leq$ \& $y_i$ \& $\leq\dots\leq$ \& $y_k$ \\
            };
            \node[AlertOrange, rotate=90] at ($(m-1-5)!0.5!(m-2-5)$) {$<$};

            \begin{scope}[on background layer]
                \node[fit=(m-1-1)(m-1-3), mysubset, "{$X_{i-1}$}" {left, UniBlue}] {};
                \node[fit=(m-2-1)(m-2-5), mysubset, AlertOrange, fill=AlertOrange!10, "{$Y_i$}" {left, AlertOrange}] {};
            \end{scope}
        \end{tikzpicture}
    \end{center}

    \newcommand{\Xmi}{\textcolor{UniBlue}{X_{i-1}}}
    \newcommand{\Yi}{\textcolor{AlertOrange}{Y_i}}
    独立集合$\Xmi$と$\Yi$に公理(I2)を適用すると，ある$y \in \Yi \setminus \Xmi$で$\Xmi + y \in \caI$となるものが存在する．
    
    \bigskip
    補題より，$y$は$x_i$を調べる時点（暫定解$X_{i-1}$）でまだ調べられていないので，$w(y) \geq w(x_i)$が成立．一方，$w(y) \leq w(y_i) \textcolor{AlertOrange}{{}<{}} w(x_i)$なので，矛盾．\qed
    % \bigskip
    % $y$を調べたときの暫定解を$X' \subseteq \Xmi$とすると，$X' + y \subseteq \Xmi + y \in \caI$であるから，公理(I1)より$X' + y \in \caI$. よって，貪欲法は$y$を$X'$に追加する．ところが，これは$y \in \Xmi$を意味するので，矛盾．\qed
\end{frame}

\begin{frame}{貪欲法の計算量}
    \small
    一般に，マトロイド$\bM = (E, \caI)$は$|E|$に対し指数個の独立集合を持つため，独立集合の一覧をアルゴリズムに入力することはできない．

    \begin{columns}
    \begin{column}{.7\textwidth}
    \begin{block}{独立性オラクル}
        部分集合$X \subseteq E$を入力すると，$X$が独立集合か否かを返してくれるオラクル
    \end{block}
    \end{column}
    \begin{column}{.29\textwidth}
        \centering
        \begin{tikzpicture}
            \node[](alg){\coloremoji[scale=2]{🖥️}};
            \node[right=2.5cm of alg](ora){\coloremoji[scale=2]{🧙‍♀️}};
            \draw[myedges, ->] (alg) -- node[above]{\footnotesize $X \in \caI$か?} (ora);
            \draw[myedges, ->] ([yshift=-5pt]ora.west) -- node[below]{\footnotesize YES/NO} ([yshift=-5pt]alg.east);
        \end{tikzpicture}
    \end{column}
    \end{columns}
    \smallskip
    具体的なマトロイド（グラフ的マトロイドや線形マトロイドなど）では，グラフの操作や行列計算により多項式時間で実装できる．

    \pause
    \bigskip
    \begin{theorem}
        貪欲法は$O(n \log n + n \mathrm{IO})$時間で最小重み独立集合を求める．ただし，$n = |E|$，$\mathrm{IO}$は独立性オラクルの計算時間．
    \end{theorem}
\end{frame}

% \begin{frame}{貪欲法によるマトロイドの特徴づけ}
%     実は「マトロイド $=$ 貪欲法で最小重み問題が解ける集合族」である．

%     \begin{definition}
%         集合族$\caI \subseteq 2^V$が\alert{独立性システム (independence system)} $\defiff$ 
%         \begin{description}
%             \item[(I0)] $\emptyset \in \caI$
%             \item[(I1)] $I \in \caI$, $I' \subseteq I$ $\implies$ $I' \in \caI$
%         \end{description}
%     \end{definition}

%     \begin{theorem}[Rado (1957), Edmonds (1971)]
%         独立性システム$\caI$に対し，以下は同値:
%         \begin{enumerate}
%             \item $\bM = (V, \caI)$はマトロイド．
%             \item 任意の重み関数$w: V \to \R$に対し，貪欲法は$\caI$の最小重み集合を返す．
%         \end{enumerate}
%     \end{theorem}
% \end{frame}

%%%%%%%%%%%%%%%%%%%%%%%%%%%%%%%%%%%%%%%%%%%%%
\section{マトロイド多面体}
\againframe<4|handout:0>{outline}
%%%%%%%%%%%%%%%%%%%%%%%%%%%%%%%%%%%%%%%%%%%%%
\begin{frame}{マトロイド多面体}
    \begin{definition}
        マトロイド$\bM$の\alert{マトロイド多面体 (matroid polytope)}を次で定める:
        \[
            P(\bM) := \conv\{\ones_I: I \in \caI \} \subseteq \R^E
        \]
        {\footnotesize \hfill ($\ones_I \in \{0,1\}^E$: $I$の特性ベクトル)}
    \end{definition}

    \bigskip
    \structure{この節の目標}
    \begin{theorem}[Edmonds (1970)]
        \[
            P(\bM) = \{x \in \R^E: x(X) \leq r_{\bM}(X) \quad (X \subseteq E), \; x \geq 0 \}
        \]
    \end{theorem}
\end{frame}

\begin{frame}{マトロイド多面体の例}
   
    % \small
    % \begin{example}[分割マトロイド]
    %     分割$(E_1, E_2, \dots, E_k)$に対する分割マトロイドのマトロイド多面体は次で表される:
    %     \begin{alignat*}{2}
    %         x(E_i \cap X) &\leq \abs{E_i \cap X} &\quad (X \subseteq E, i = 1, \dots, k)\\
    %         x &\geq 0 &
    %     \end{alignat*}
    % \end{example}
    
    \begin{example}[グラフ的マトロイド]
        連結な無向グラフ$G = (V, E)$に対するグラフ的マトロイドのマトロイド多面体は次で表される:
        \begin{alignat*}{2}
            x(X) &\leq \abs{V(G[X])} - c(G[X]) &\quad (X \subseteq E)\\
            x &\geq 0 &
        \end{alignat*}
        {\footnotesize \hfill ($G[X]$: $X$の誘導部分グラフ, $c(G[X])$: $G[X]$の連結成分数)}
    \end{example}
\end{frame}

\begin{frame}{定理の証明 (1/4)}
    $P := \{x \in \R^E: x(X) \leq r_{\bM}(X) \quad (X \subseteq E), \quad x \geq 0 \}$とする．

    \bigskip
    \structure{$P(\bM) \subseteq P$であること:}
    $I \in \caI$を独立集合とし，$x = \ones_I$を特性ベクトルとする．任意の部分集合$X \subseteq V$に対し，
    \begin{align*}
        x(X) = |I \cap X| \leq r_{\bM}(X). \tag{$I \cap X \subseteq X$は独立集合より}
    \end{align*}
    ∴$x \in P$. $P$は凸多面体なので，$P(\bM) \subseteq P$.
\end{frame}

\begin{frame}{定理の証明 (2/4)}
    \small
    \structure{$P \subseteq P(\bM)$であること:} 任意の$w \in \R^E$に対し，$\max_{x \in P} w^\top x$を達成する点が$P(\bM)$に存在することを示す．

    \bigskip
    $w$を重みとする最大重み独立集合問題を考える:
    \begin{block}{}
        \[
            \text{maximize} \quad w(I) \quad\text{s.t.}\quad I \in \caI 
        \]
    \end{block}
   
    \bigskip
    重みを降順に並べて，
    \[  
        w(e_1) \geq w(e_2) \geq \dots \geq w(e_k) \geq 0 > w(e_{k+1}) \geq \dots \geq w(e_n)
    \]
    とすると，$e_1, \dots, e_k$まで貪欲法を走らせた暫定解が最適解である．それを
    \[
        I = \{e_{i_1}, e_{i_2}, \dots, e_{i_\ell}\} \quad (1 \leq i_1 < i_2 < \dots < i_\ell \leq k) 
    \]
    とおく．$x = \ones_I \in P(\bM)$が$\max_{x \in P} w^\top x$の最適解であることを示せば良い．
\end{frame}

\begin{frame}{定理の証明 (3/4)}
    \footnotesize
    \begin{primal}
        \setlength{\abovedisplayskip}{0pt}
        \begin{alignat*}{3}
            \text{max}   \quad & w^\top x && \quad \\
            \text{s.t.}  \quad & x(X)  \leq r(X) && \quad (X \subseteq E) \\
                         \quad & x(e) \geq 0 && \quad (e \in E)
        \end{alignat*}
    \end{primal}
    \begin{dual}
        \setlength{\abovedisplayskip}{0pt}
        \begin{alignat*}{3}
            \text{min}  \quad & \sum_{X \subseteq E} r(X) y_X && \\
            \text{s.t.} \quad & \sum_{X \ni e} y_X \geq w(e)  &&\quad && (e \in E) \\
                         \quad & y_X  \geq 0                  &&\quad && (X \subseteq E)
        \end{alignat*}
    \end{dual}
\end{frame}

\begin{frame}{定理の証明 (4/4)}
    \small
    LPの一般論より，
    $x = \ones_I$が主最適解であることを示すには，強双対性$w^\top x = \sum_{X \subseteq E} r(X)y_X$を満たす双対実行可能解$y$を構成すればよい．

    \bigskip
    実際，
    \begin{align*}
        y_S := \begin{cases}
        w(e_{i_j}) - w(e_{i_{j+1}}) & \text{if } S = \{e_{i_1}, \dots, e_{i_j}\} \quad (j = 1, \dots, \ell - 1) \\
        w(e_{i_\ell})                 & \text{if } S = \{e_{i_1}, \dots, e_{i_\ell}\} \\
        0                             & \text{otherwise}
        \end{cases}
    \end{align*}
    とすると，$y$は双対実行可能で，$(x, y)$が強双対性を満たすことが，直接計算により確認できる（各自確認せよ）． \qed
\end{frame}

\begin{frame}{マトロイド多面体上の線形最適化}
    \begin{corollary}
        マトロイド多面体上の線形最適化
        \[  
              \max \quad w^\top x \quad\text{s.t.}\quad x(X) \leq r_{\bM}(X) \quad (X \subseteq E), \; x \geq 0
        \] 
        は$O(n \log n + n \mathrm{IO})$時間で最適解を求められる．
    \end{corollary}

    \bigskip
    指数本の不等式があるにもかかわらず，貪欲法で最適解を求められる．
\end{frame}

\begin{frame}{マトロイド多面体の整数性}
    \begin{corollary}
        不等式系$\{x(X) \leq r_{\bM}(X) \quad (X \subseteq E), \; x \geq 0\}$は整数多面体を定める．
    \end{corollary}

    \bigskip
    $P(\bM)$を定める不等式系の係数行列は\alert{完全単模行列ではない！}\\
    {\footnotesize \structure{(∵)}たとえば$\begin{bmatrix}
        1 & 1 & 0 \\
        0 & 1 & 1 \\
        1 & 0 & 1
    \end{bmatrix}$のように，行列式$=2$の部分行列を含んでいる．}

    \bigskip
    マトロイド多面体は完全単模性に依らず，整数性が証明できる．
\end{frame}

\begin{frame}{マトロイド基多面体}
    基についても同様の定理が成り立つ．

    \begin{definition}
        マトロイド$\bM$の\alert{基多面体 (base polytope)}を次で定める:
        \[
            B(\bM) := \conv\{\ones_B: B \in \caB \} \subseteq \R^E
        \]
        {\footnotesize \hfill ($\ones_B \in \{0,1\}^E$: $B$の特性ベクトル)}
    \end{definition}

    \bigskip
    \begin{theorem}[Edmonds (1970)]
        \setlength{\abovedisplayskip}{0pt}
        \begin{align*}
            B(\bM) = \{x \in \R^E: x(X) \leq r_{\bM}(X) \; (X \subseteq E), \; \highlight{$x(E) = r_{\bM}(E)$}, \; x \geq 0 \}
        \end{align*}
    \end{theorem}
\end{frame}

\againframe<1>{outline}
\end{document}